\documentclass{article}
\usepackage[UTF8]{ctex}
\usepackage{amsmath,amssymb,amsfonts}
\usepackage{enumitem}
\usepackage{geometry}
\geometry{a4paper, margin=1in}
% === 错题本排版优化（不参与题目标识） ===
\setlength{\parindent}{0pt}
\setlength{\parskip}{0.35em plus 0.10em minus 0.08em}
\setlength{\emergencystretch}{3em}
\setlist{itemsep=0.25em, topsep=0.40em, parsep=0pt}
\allowdisplaybreaks[2]
% === 错题本排版优化结束 ===
\begin{document}
\section*{四川大学2026年数学分析真题}
\section*{166．四川大学2026年数学分析试卷}

\subsection{题目1}
1．（16 分）求极限：

\begin{itemize}
  \item[（1）] \(\lim _{n \rightarrow \infty}\left(\frac{1}{\ln (n+\pi)-\ln n}-\frac{n}{\pi}\right)\) ．
  \item[（2）] \(\lim _{x \rightarrow 0^{+}} \frac{\int_{0}^{x^{2}}(x-t)(1-\cos \sqrt{t}) \mathrm{d} t}{\int_{0}^{x} \tan ^{4} t \sec t \mathrm{~d} t}\) ．
\end{itemize}

\subsection{题目2}
2．（12 分）计算 \(I_{n}=\int_{0}^{\pi} \cos ^{n} x \mathrm{~d} x\) ．

\subsection{题目3}
3．（12 分）要做一个容器为 \(1 \mathrm{~m}^{3}\) 的无盖铝圆桶，什么样的尺寸才能使用料最省？

\subsection{题目4}
4．（12 分）计算

\[
\int_{L}(x-y) \mathrm{d} x+z^{3} \mathrm{~d} z
\]

其中 \(L\) 是 \(x^{2}+y^{2}+z^{2}=1\) 与 \(z=c(|c|<1)\) 的交线，从 \(z\) 轴的正向往负向看去是逆时针方向。

\subsection{题目5}
5．（12 分）计算

\[
\iint_{\Sigma} \frac{x \mathrm{~d} y \mathrm{~d} z+y \mathrm{~d} z \mathrm{~d} x+z \mathrm{~d} x \mathrm{~d} y}{\left(x^{2}+y^{2}+z^{2}\right)^{\frac{3}{2}}} .
\]

其中 \(\sum: 1-\frac{z}{3}=\frac{(x-2)^{2}}{25}+\frac{(y-1)^{2}}{16}(z \geq 0)\) ，方向取上侧。

\subsection{题目6}
6．（13 分）求幂级数 \(\sum_{n=0}^{\infty}(-1)^{n} \frac{x^{n+3}}{(n+1)(n+3)}\) 的收敛域与和函数．

\subsection{题目7}
7．（13 分）判断积分 \(\int_{1}^{+\infty} x^{3} \sin \left(x^{q}\right) \mathrm{d} x(q \neq 0)\) 的敛散性（包括绝对收敛，条件收敛和发散）．

\subsection{题目8}
8．（12 分）证明：方程 \(e^{x}=x^{2}+5 x+1\) 的实根不超过3个。

\subsection{题目9}
9．（12 分）设 \(f:[0,+\infty) \rightarrow[0,+\infty)\) 连续，且 \(\lim _{x \rightarrow+\infty} f(x) \int_{0}^{x} f(t) \mathrm{d} t=a \in(0,+\infty)\) ，证明： \(\lim _{x \rightarrow+\infty} \sqrt{x} f(x)\) 存在且为正。

\subsection{题目10}
10．（12 分）设 \(f(x)\) 在 \([0,+\infty)\) 上一致连续，且对任意固定的 \(x \in[0,+\infty)\) ，当自然数 \(n \rightarrow+\infty\) 时，有\(f(x+n) \rightarrow 0\) ．证明：函数列 \(\{f(x+n)\}_{n=1}^{\infty}\) 在 \([0,1]\) 上一致收敛于 0 ．

\subsection{题目11}
11．（14 分）设 \(f(x, y)\) 在 \(\left(x_{0}, y_{0}\right)\) 处可微，且 \(f\left(x_{0}, y_{0}\right)=0 . g(x, y)\) 在 \(\left(x_{0}, y_{0}\right)\) 处连续．证明：\(f(x, y) g(x, y)\) 在 \(\left(x_{0}, y_{0}\right)\) 处可微，且 \(\mathrm{d}(f g)\left(x_{0}, y_{0}\right)=g\left(x_{0}, y_{0}\right) \mathrm{d} f\left(x_{0}, y_{0}\right)\) ．

\subsection{题目12}
12．（10 分）设 \(f(x)\) 为［0，1］上非负且严格单调增加的连续函数。证明：对任意的正整数 \(n\) ，存在唯一的 \(x_{n} \in[0,1]\) ，使得 \(\left(f\left(x_{n}\right)\right)^{n}=\int_{0}^{1}(f(x))^{n} \mathrm{~d} x\) ，且 \(\lim _{n \rightarrow \infty} x_{n}=1\) ．
\end{document}
